% !TEX root = ../notes.tex
\documentclass[11pt]{article}

\usepackage[a4paper,margin=1in]{geometry}
\usepackage{amsmath,amssymb,amsfonts,amsthm}
\usepackage{mathtools}
\usepackage{enumitem}
\usepackage{microtype}
\usepackage{bm}
\usepackage{booktabs}
\usepackage{xcolor}
\usepackage{hyperref}
\hypersetup{colorlinks=true,linkcolor=black,citecolor=black,urlcolor=black}

\usepackage{preamble}

\setlist[itemize]{topsep=2pt,itemsep=2pt,parsep=0pt,leftmargin=1.2em}
\setlist[enumerate]{topsep=2pt,itemsep=2pt,parsep=0pt,leftmargin=1.4em}

\title{\vspace{-0.7em}\Huge Demographic Parity in regression \& classiffication (unawareness)\\
\large{\emph{}}}
\date{}

\begin{document}
\maketitle

\part{Framework from \textcite{divol2024demographicparityregressionclassification}} 
\section{Setting}
Data $(X,S,Y)\in\X\times\Sset\times\R$; $S\in\{1,\dots,K\}$.
\begin{itemize}
  \item Group weights $p_s:=\Pp(S=s)$, feature laws $\mu_s:=Law(X\mid S=s)$, overall $\mu=\sum_s p_s\,\mu_s$.
  \item Bayes regressor (unconstrained): $\eta(x):=\E[Y\mid X=x]$.
  \item For $f:\X\to\R$, squared risk $R(f)=\E[(Y-f(X))^2] = \E[(Y-\eta(X))^2]+\E[(\eta(X)-f(X))^2]$.
  \item Pushforward: $f_{\push}\mu$ is law of $f(X)$ when $X\sim\mu$.
\end{itemize}

\begin{figure}[h]
  \centering
    \centering
    \includegraphics[width=0.7\linewidth]{figs/paper/OT/rf/10000/example_base_pred_dists.pdf}
    \caption{Distributions of $\eta(X)$ in each group (colors) and overall (black).}
\end{figure}

\begin{figure}[h]
  \centering
    \centering
    \includegraphics[width=\linewidth]{figs/paper/OT/rf/10000/example_scatter_XY_colored_by_S.pdf}
    \caption{Scatter Y as a function of each column of X (left) on train (right) on test.}
\end{figure}

\begin{figure}
  \centering
    \centering
    \includegraphics[width=\linewidth]{figs/paper/confounded/rf/scatter_XY_by_gamma/scatter_XY_gamma0p00_beta1.0.pdf}
    \caption{Scatter Y as a function of X ($\gamma=0$ ) .}
\end{figure}

\begin{figure}
  \centering
    \centering
    \includegraphics[width=\linewidth]{figs/paper/confounded/rf/scatter_XY_by_gamma/scatter_XY_gamma5p00_beta1.0.pdf}
    \caption{Scatter Y as a function of X ($\gamma=5$ ) .}
\end{figure}


\section{DP}
\textbf{DP for regression (unawareness):} $f(X)\perp S$,
i.e.\ $f_{\push}\mu_s$ is the same for all $s$.
\begin{equation}
\min_{f:\X\to\R}\ \E[(Y-f(X))^2]\quad\text{s.t.}\quad f(X)\perp S.
\end{equation}
\emph{Order preservation (within group):} if $\eta(X)<\eta(X')$ and $S=S'$, then $f(X)<f(X')$ a.s.

\medskip
\textit{Key impossibility:} if $f$ has DP while $\eta$ is not DP, then $f$ cannot preserve order (proof by contrapositive: DP+$\uparrow$ order $\Rightarrow$ $\eta$ DP).

\section{Two groups $\Rightarrow$ reduce to ``excess-mass'' pieces}
Let $\alpha:=\mu_1-\mu_2$ (\emph{signed} measure). Jordan decomposition $\alpha=\alpha^+-\alpha^-$ with equal mass $m:=\alpha^+(\X)=\alpha^-(\X)=\tfrac12\|\mu_1-\mu_2\|_{\TV}$.
\begin{equation}
\mu_+:=\alpha^+/m,\qquad \mu_-:=\alpha^-/m.
\end{equation}
Supports split $\X$ into $\X_+$ (where group 1 dominates), $\X_-$ (group 2 dominates), and the tie region $\X_=$.
\begin{equation}
\boxed{~f\ \text{is DP} \iff (f_{\push}\mu_+)=(f_{\push}\mu_-)~}
\end{equation}
(\emph{So DP is governed by how $f$ acts on the disagreement regions; $\X_=$ is irrelevant to DP}.)

\section*{3. The 2D lift and the OT cost (this is the main trick)}
Pick a reference $\mu$ dominating $\mu_\pm$. Define the \emph{signed density gap}
\begin{equation}
\Delta(x):=\begin{cases}
\dfrac{\dd\mu_+}{\dd\mu}(x), & x\in\X_+,\\[4pt]
-\dfrac{\dd\mu_-}{\dd\mu}(x), & x\in\X_-,\\
0, & x\in\X_=,
\end{cases}
\qquad\text{and note}\quad
\dd\mu=\frac{1}{|\Delta|}\dd\mu_\pm\ \text{on }\X_\pm.
\end{equation}
Lift $x\mapsto x^{\uparrow}:=(\eta(x),\Delta(x))=(h,d)\in\Omega:=\{(h,d)\in\R^2:\,d\neq 0\}$.
Let $\mu_\pm^{\uparrow}$ be the laws of $x^\uparrow$ under $\mu_\pm$.

\begin{figure}[h]
    \centering
    \includegraphics[width=\linewidth]{figs/paper/OT/rf/10000/example_eta_delta_scatter.pdf}
  \caption{Scatter of $(\eta,\Delta)$ on train (left) and test (right).}
\end{figure}

\medskip
\textbf{Cost (multi$\to$one):}\quad $c((h,d),y)=\dfrac{(h-y)^2}{|d|}$.

\section*{4. Barycenter program }
\begin{equation}
\boxed{~\min_{\nu\in\mathcal P(\R)}\ \OT_c(\mu_+^{\uparrow},\nu)+\OT_c(\mu_-^{\uparrow},\nu)~}
\end{equation}
This equals the optimal \emph{excess} risk among DP regressors (first reduce risk to $\E[(\eta-f)^2]$, then use $\dd\mu=\dd\mu_\pm/|\Delta|$ on $\X_\pm$).

\section{Collapse to a standard 2-marginal OT}
For fixed $(x_1,x_2)$ with $x_i=(h_i,d_i)$,
\begin{equation}
\min_y\ \frac{(h_1-y)^2}{|d_1|}+\frac{(h_2-y)^2}{|d_2|}
\quad\Rightarrow\quad
y^\star=\frac{\frac{h_1}{|d_1|}+\frac{h_2}{|d_2|}}{\frac{1}{|d_1|}+\frac{1}{|d_2|}},
\end{equation}
and minimum value
\begin{equation}
\boxed{~C((h_1,d_1),(h_2,d_2))=\frac{(h_2-h_1)^2}{|d_1|+|d_2|}~}.
\end{equation}
Hence
\begin{equation}
\min_{\nu}\ \OT_c(\mu_+^{\uparrow},\nu)+\OT_c(\mu_-^{\uparrow},\nu)
\;=\;\OT_{C}(\mu_+^{\uparrow},\mu_-^{\uparrow}).
\end{equation}

\begin{figure}[h]
  \centering
    \centering
    \includegraphics[width=\linewidth]{figs/paper/OT/rf/10000/transport/example_ot_matrices_full.pdf}
    \caption{OT matrix between $\mu_+^{\uparrow}$ and $\mu_-^{\uparrow}$.}
\end{figure}

\begin{figure}[h]
  \centering
    \centering
    \includegraphics[width=0.7\linewidth]{figs/paper/OT/rf/10000/example_agg_plan_eta.pdf}
    \caption{Aggregated OT plan projected on the $\eta$-axis.}
    \end{figure}

    \begin{figure}[h]
      \centering
      \begin{minipage}[t]{0.49\linewidth}
        \centering
        \includegraphics[width=\linewidth]{figs/paper/OT/rf/10000/example_eta_delta_transport_lines.pdf}
        \small OT transport plan from $\mu_+^{\uparrow}$ to $\mu_-^{\uparrow}$ (top 30\% largest moves; max 150 links).
      \end{minipage}\hfill
      \begin{minipage}[t]{0.49\linewidth}
        \centering
        \includegraphics[width=\linewidth]{figs/paper/OT/rf/10000/example_barycenter_links.pdf}
        \small Barycenter links from $(\eta,\Delta)$ to the fair prediction (top 30\% largest moves; max 150 links).
      \end{minipage}
      \caption{OT transport plan (left) VS barycenter links (right).}
    \end{figure}

  \begin{figure}[h]
    \centering
    \begin{minipage}[t]{0.49\linewidth}
      \centering
      \includegraphics[width=\linewidth]{figs/paper/OT/rf/10000/example_eta_delta_colored_by_Y.pdf}
      \caption{Train on the left test on the right}
    \end{minipage}\hfill
    \begin{minipage}[t]{0.49\linewidth}
      \centering
      \includegraphics[width=\linewidth]{figs/paper/OT/rf/10000/example_pred_dists.pdf}
      \caption{Distributions before and after correction in each group.}
    \end{minipage}
  \end{figure}



\section{Existence, uniqueness, and maps (the mild regularity)}
\textbf{Assumption (zero mass on graphs).} $\mu_\pm^{\uparrow}$ give zero measure to any set $\{(h,d):h=F(d),\,d\neq 0\}$.
(this reduces to non-atomicity assumption in the aware setting)\\

\smallskip
Under this and $\E[Y^2]<\infty$:
\begin{itemize}
  \item The barycenter is \emph{unique}: $\nu^{\mathrm{bar}}$.
  \item Each $\OT_c(\mu_\pm^{\uparrow},\nu^{\mathrm{bar}})$ is realized by a \emph{transport map} $f_\pm:\Omega\to\R$.
\end{itemize}
Let $\varphi$ be a $C$-convex Kantorovich potential for $\OT_C(\mu_+^{\uparrow},\mu_-^{\uparrow})$.
Where $\partial_h\varphi(h,d)$ exists,
\begin{equation}
  \boxed{~f_+(h,d)=h-\frac{|d|}{2}\,\partial_h\varphi(h,d)~}\qquad
(\text{and symmetrically } f_-).
\end{equation}


\section*{7. Final form of the optimal fair regressor}
\begin{equation}
  f^\star(x)=
\begin{cases}
f_+\big(\eta(x),\Delta(x)\big), & x\in\X_+,\\[4pt]
f_-\big(\eta(x),\Delta(x)\big), & x\in\X_-,\\[4pt]
\eta(x), & x\in\X_=.
\end{cases}
\qquad\text{and}\qquad
(f^\star)_{\push}\mu_+=(f^\star)_{\push}\mu_-=\nu^{\mathrm{bar}}.
\end{equation}


\section{ Nestedness (multi to one OT)}
Setup: source $(X,\mu)\subset\R^m$, target $(Y,\nu)\subset\R$, cost $c(x,y)$ differentiable in $y$.
Define sublevel sets using the \emph{marginal cost}
\begin{equation}
X_{\le}(y,k):=\{x\in\X:\,\partial_y c(x,y)\le k\}.
\end{equation}
For each $y$, choose $k(y)$ so that $\mu\big(X_{\le}(y,k(y))\big)=F_\nu(y)$ where $F_\nu$ is the CDF of $\nu$.

\textbf{Nestedness:}
\begin{equation}
y_1<y_2\ \Longrightarrow\ X_{\le}\!\big(y_1,k(y_1)\big)\subseteq X_{\le}\!\big(y_2,k(y_2)\big).
\end{equation}
Consequence: existence of a Monge map $T:\X\to Y$ with
$\partial_y c\big(x,T(x)\big)=k\big(T(x)\big)$ a.e.\ (no mass splitting).

\medskip
\emph{Tiny example.} If $c(x,y)=\tfrac12\|x\|^2-y\,x_1$, then $\partial_y c(x,y)=-x_1$, so $X_{\le}(y,k)=\{x_1\ge -k\}$ (independent of $y$). Picking $k(y)$ to match $F_\nu(y)$ makes these sets nested $\Rightarrow$ a map depending on $x_1$ only.


\begin{table}[t]
  \centering
  \begin{tabular}{lrrrr}
  \toprule
  Subset & Spearman $\rho$ & Kendall $\tau_b$ & Mean $|\Delta\mathrm{rank}|$ & Top-5\% Jaccard \\
  \midrule
  overall & 0.662 & 0.509 & 0.216 & 0.250 \\
  S=0 & 0.980 & 0.883 & 0.044 & 0.682 \\
  S=1 & 0.979 & 0.881 & 0.045 & 0.805 \\
  \bottomrule
  \end{tabular}
  \caption[Ordering stability between eta-hat and corrected scores.]{Ordering stability between $\hat\eta$ and corrected scores.}
  \label{tab:ordering}
\end{table}

\begin{table}[t]
  \centering
  \begin{tabular}{lr}
  \toprule
  Quantity & AUC \\
  \midrule
  $\mathrm{AUC}(S\mid Y)$ & 0.858 \\
  $\mathrm{AUC}(S\mid \hat{\eta})$ & 0.828 \\
  $\mathrm{AUC}(S\mid \tilde{y})$ & 0.396 \\
  \bottomrule
  \end{tabular}
  \caption[Leakage AUC summary.]{How easy it is to infer S from te score}
  \label{tab:auc}
\end{table}

  $\implies$ mélange entre les groupes mais dans chaque groupe l'ordre est assez preservé 

  
\section{Classification link (cost-sensitive, $Y\in\{0,1\}$)}
For weight $y\in[0,1]$,
\begin{equation}
R_y(g)=y\,\Pp(Y=0,\,g(X)=1)+(1-y)\,\Pp(Y=1,\,g(X)=0),\quad g:\X\to\{0,1\}.
\end{equation}
Under DP and a \emph{nestedness} condition, the optimal fair classifiers are \emph{thresholds} of $f^\star$:
\begin{equation}
g_y^\star(x)=\1\{\,f^\star(x)\ge t(y)\,\},
\end{equation}
and the family $\{g_y^\star\}_y$ recovers $f^\star$ (up to null sets). Without nestedness this equivalence can fail.

\section{Awareness VS unawareness }
\begin{itemize}
  \item \textbf{Awareness} (can use $S$): DP-optimal regressor is a \emph{groupwise} OT/quantile match to a common barycenter; preserves within-group order; often closed-form via CDF$\leftrightarrow$quantile.
  \item \textbf{Unawareness}: DP forces a \emph{2D$\to$1D} OT barycenter depending on $(\eta,\Delta)$; no general closed form; order preservation can \emph{fail}.
\end{itemize}

\section{ In practice }
\begin{enumerate}
  \item Estimate $\widehat\eta(x)$ and $\widehat\Delta(x)$ (need $\widehat{\Pp}(S=s\mid X=x)$ or density ratios).
  \item Build empirical $\widehat\mu_\pm^{\uparrow}$ on samples from $\X_\pm$.
  \item Solve $\OT_C$ with $C((h_1,d_1),(h_2,d_2))={(h_2-h_1)^2}/{(|d_1|+|d_2|)}$ to get maps $\hat f_\pm$.
  \item Output $\hat f^\star$ piecewise on $(\X_+,\X_-,\X_=)$ as above.
\end{enumerate}


\section{How good is the method when the X by S are more or less confounded? }

All confounded experiements are with $n=1500$

\begin{figure}[h]
  \centering
    \centering
    \includegraphics[width=\linewidth]{figs/paper/confounded/rf/sweep_OT_n1500_beta1.0_100g.pdf}
    \caption{Results for different gamma values (make sense bc when $\gamma=0$ fairness is free).}
\end{figure}

\begin{figure}[h]
  \centering
  \begin{minipage}[t]{0.51\linewidth}
    \centering
    \includegraphics[width=1.2\linewidth]{figs/paper/confounded/rf/scatter_XY_by_gamma/scatter_XY_gamma0p00_beta1.0.pdf}
  \end{minipage}\hspace{0.01cm}
  \begin{minipage}[t]{0.48\linewidth}
    \centering
    \includegraphics[width=0.61\linewidth]{figs/paper/confounded/rf/eta_delta_by_gamma/eta_delta_gamma0p00_beta1.0.pdf}
  \end{minipage}
  \caption{Results at $\gamma=0$: scatter $(X,Y)$ vs. $(\eta,\Delta)$ side by side.}
\end{figure}

\begin{figure}[h]
  \centering
  \begin{minipage}[t]{0.51\linewidth}
    \centering
    \includegraphics[width=1.2\linewidth]{figs/paper/confounded/rf/scatter_XY_by_gamma/scatter_XY_gamma1p75_beta1.0.pdf}

  \end{minipage}\hspace{0.01cm}
  \begin{minipage}[t]{0.48\linewidth}
    \centering
    \includegraphics[width=0.61\linewidth]{figs/paper/confounded/rf/eta_delta_by_gamma/eta_delta_gamma1p75_beta1.0.pdf}
  \end{minipage}
  \caption{Results at $\gamma=1.88$: scatter $(X,Y)$ vs. $(\eta,\Delta)$ side by side.}
\end{figure}


\begin{figure}[h]
  \centering
  \begin{minipage}[t]{0.51\linewidth}
    \centering
    \includegraphics[width=1.2\linewidth]{figs/paper/confounded/rf/scatter_XY_by_gamma/scatter_XY_gamma2p78_beta1.0.pdf}

  \end{minipage}\hspace{0.01cm}
  \begin{minipage}[t]{0.48\linewidth}
    \centering
    \includegraphics[width=0.61\linewidth]{figs/paper/confounded/rf/eta_delta_by_gamma/eta_delta_gamma2p78_beta1.0.pdf}
  \end{minipage}
  \caption{Results at $\gamma=2.50$: scatter $(X,Y)$ vs. $(\eta,\Delta)$ side by side.}
\end{figure}


\begin{figure}[h]
  \centering
  \begin{minipage}[t]{0.51\linewidth}
    \centering
    \includegraphics[width=1.2\linewidth]{figs/paper/confounded/rf/scatter_XY_by_gamma/scatter_XY_gamma3p89_beta1.0.pdf}

  \end{minipage}\hspace{0.01cm}
  \begin{minipage}[t]{0.48\linewidth}
    \centering
    \includegraphics[width=0.61\linewidth]{figs/paper/confounded/rf/eta_delta_by_gamma/eta_delta_gamma3p89_beta1.0.pdf}
  \end{minipage}
  \caption{Results at $\gamma=3.75$: scatter $(X,Y)$ vs. $(\eta,\Delta)$ side by side.}
\end{figure}


\begin{figure}[h]
  \centering
  \begin{minipage}[t]{0.51\linewidth}
    \centering
    \includegraphics[width=1.2\linewidth]{figs/paper/confounded/rf/scatter_XY_by_gamma/scatter_XY_gamma5p00_beta1.0.pdf}

  \end{minipage}\hspace{0.01cm}
  \begin{minipage}[t]{0.48\linewidth}
    \centering
    \includegraphics[width=0.61\linewidth]{figs/paper/confounded/rf/eta_delta_by_gamma/eta_delta_gamma5p00_beta1.0.pdf}
  \end{minipage}
  \caption{Results at $\gamma=5.00$: scatter $(X,Y)$ vs. $(\eta,\Delta)$ side by side.}
\end{figure}

\section{defs}
\begin{itemize}
  \item \textbf{DP:} $f(X)\perp S$ $\iff$ $f_{\push}\mu_s$ independent of $s$.
  \item \textbf{Pushforward:} $f_{\push}\mu(A)=\mu(f^{-1}(A))$.
  \item \textbf{Jordan decomposition:} $\alpha=\alpha^--\alpha^+$ with disjoint supports, $|\alpha|=\alpha^++\alpha^-$.
  \item \textbf{Barycenter (here):} $\nu$ minimizing $\OT$-distance sum from sources.
  \item \textbf{Nestedness:} monotone growth in $y$ of $\partial_y c$-sublevel sets with prescribed $\mu$-mass.
  \item \textbf{Kantorovich potential:} dual function whose subgradient characterizes optimal plans/maps.
\end{itemize}

\section{formuless}
\begin{equation}
\boxed{~c((h,d),y)=\frac{(h-y)^2}{|d|}~},\qquad
\boxed{~C((h_1,d_1),(h_2,d_2))=\frac{(h_2-h_1)^2}{|d_1|+|d_2|}~},
\end{equation}
\begin{equation}
\boxed{~f^\star(x)=
\begin{cases}
\eta(x)-\dfrac{|\Delta(x)|}{2}\,\partial_h\varphi(\eta(x),\Delta(x)), & x\in\X_+\cup\X_-,\\[4pt]
\eta(x), & x\in\X_=.
\end{cases}~}.
\end{equation}

\newpage


\clearpage
\part{Unawareness from \textcite{chzhen2024regressiondemographicparityconstraints}}
\label{part:chzhen}



\section{what porblem do they solve?}
We want regression under \textbf{demographic parity (DP)} in the \textbf{unawareness} setting:
the predictor cannot use $S$ at inference time, but we still want $\widehat Y \perp S$.

Key design choice: build a \emph{smooth convex dual objective} whose stochastic gradients can be
computed from \emph{unlabeled} $X$ only, enabling \textbf{online post-processing} on a stream of i.i.d.\ $X$.

\section{Setting and notation}
Data $(X,S,Y)\in \R^d \times [K]\times \R$ with unknown law $P$.
\begin{itemize}
  \item Regression function: $\eta(x) := \E[Y\mid X=x]$.
  \item Sensitive marginals: $p_s := \Pp(S=s)$, and vector $p=(p_s)_{s\in[K]}$.
  \item Sensitive predictor / propensity: $\tau(x) := (\tau_s(x))_{s\in[K]}$ with $\tau_s(x):=\Pp(S=s\mid X=x)$.
\end{itemize}

\paragraph{Randomized predictors (unawareness).}
A randomized prediction function is a Markov kernel
\begin{equation}
\pi: 
\begin{cases}
  \mathcal{B}(\R)\times \R^d &\to [0,1],\\
  \qquad B&\mapsto \pi(B\mid x).
\end{cases}
\end{equation}
It induces a prediction random variable $\widehat Y_\pi$ by
\begin{equation}
Law(\widehat Y_\pi \mid X=x, S=s)=\pi(\cdot\mid x).
\end{equation}
\textbf{Remark 2.1:} deterministic predictors correspond to Dirac kernels
$\pi(\cdot\mid x)=\delta_{g(x)}$, giving $\widehat Y_\pi=g(X)$. Also, $\pi$ specifies the joint law
of $(X,S,\widehat Y_\pi)$, and to relate it to $(X,S,Y)$ one may assume conditional independence
$\widehat Y_\pi \perp Y \mid (X,S)$ so that $(X,S,Y,\widehat Y_\pi)$ is well-defined.

\section{Demographic parity and risk}
\paragraph{DP.}
A predictor $\pi$ satisfies demographic parity if
\begin{equation}
\widehat Y_\pi \perp S.
\end{equation}
This is \emph{distributional}, not pointwise in $x$.

\paragraph{Risk.}
They measure performance against $\eta$ (excess risk):
\begin{equation}
R(\pi) := \E\big[(\widehat Y_\pi-\eta(X))^2\big]
        = \E\Big[\int_{\R} (\widehat y-\eta(X))^2\,\pi(d\widehat y\mid X)\Big].
\end{equation}

\section{discretization of the prediction range}
They force the predictor outputs to lie on a \textbf{fixed finite grid}.

\paragraph{Bounded signal assumption.}
Assume there exists $B>0$ such that $|\eta(X)|\le B$ a.s.

\paragraph{Grid.}
For an integer $L\ge 1$, define the uniform grid on $[-B,B]$
\begin{equation}
\widehat{\mathcal Y}_L := \left\{ \widehat y_\ell := \frac{\ell B}{L} : \ell\in [[L]]:=\{-L,\ldots,L\}\right\},
\qquad |\widehat{\mathcal Y}_L|=2L+1.
\end{equation}
A feasible predictor must satisfy $Supp(\pi(\cdot\mid x))=\widehat{\mathcal Y}_L$ for all $x$.

\paragraph{Why discretize?}
Because with a fixed finite support, DP reduces to finitely many \emph{linear} constraints on the
probabilities $\pi(\ell\mid x)$:
\begin{equation}
\Pp(\widehat Y_\pi=\widehat y_\ell\mid S=s)=\Pp(\widehat Y_\pi=\widehat y_\ell)
\quad\Longleftrightarrow\quad
\E[\pi(\ell\mid X)\mid S=s]=\E[\pi(\ell\mid X)].
\end{equation}
They allow slack via thresholds $\varepsilon=(\varepsilon_s)_{s\in[K]}$ and define (per group, per bin)
\begin{equation}
U_s(\pi,\ell):=\Big|\E[\pi(\ell\mid X)\mid S=s]-\E[\pi(\ell\mid X)]\Big| \le \varepsilon_s.
\end{equation}

\paragraph{Important sanity point \marie{(parce que j'ai eu un doute)}: this is \emph{not} turning regression into classification.}
The algorithm learns a \emph{distribution} over grid values; you can output either a sample
$\widehat Y_\pi$ or the mean $\E[\widehat Y_\pi\mid X=x]=\sum_{\ell} \widehat y_\ell\,\pi(\ell\mid x)$,
which is still a real-valued regressor. The discretization is an algorithmic device to make DP constraints finite.

\section{Entropic regularization $\Rightarrow$ smooth dual convex program}
Solving the discretized constrained problem directly is a large LP in dimension $\Omega(LK)$, so they
switch to an entropic-regularized variant \marie{(comme c'est leur gros argument je pense que dans la comparaison il faudra voir la complexité temporelle \emph{a minima})}.

\paragraph{Entropy-regularized risk.}
Let $\Psi(\mu)=\sum_{y\in Supp(\mu)}\mu(y)\log\mu(y)$ denote negative entropy. For $\beta>0$:
\begin{equation}
R_\beta(\pi):=R(\pi) + \frac{1}{\beta}\E\big[\Psi(\pi(\cdot\mid X))\big].
\end{equation}

\paragraph{Dual variables and key reparameterization.}
Define
\begin{equation}
t(x):= \mathbf{1} - \frac{\tau(x)}{p}\in\R^K \qquad (\text{componentwise division}),
\qquad
r_\ell(x):=\big(\eta(x)-\widehat y_\ell\big)^2.
\end{equation}

\subsection{Lemma 3.1 (the main structural lemma)}
The entropic-regularized problem admits a solution that can be written explicitly in terms of
nonnegative dual matrices $\Lambda=(\lambda_{\ell s})$ and $V=(\nu_{\ell s})$, as follows. \marie{2 matrices parce que contrainte en valeur absolue}

\paragraph{Dual objective (smooth convex).}
Let $\lambda_\ell=(\lambda_{\ell s})_{s\in[K]}$ and $\nu_\ell=(\nu_{\ell s})_{s\in[K]}$.
Define the log-sum-exp at temperature $\beta$ by
\begin{equation}
\mathrm{LSE}_\beta(w)=\beta^{-1}\log\Big(\sum_j e^{\beta w_j}\Big),
\end{equation}
and the soft-argmax / softmax by $\sigma_j(w)=\frac{e^{w_j}}{\sum_i e^{w_i}}$.
Then define
\begin{equation}
F(\Lambda,V):=
\E\Big[
\mathrm{LSE}_\beta\Big( \big(\langle \lambda_\ell-\nu_\ell,\,t(X)\rangle - r_\ell(X)\big)_{\ell\in[[L]]}\Big)
\Big]
\;+\;
\sum_{\ell\in[[L]]}\langle \lambda_\ell+\nu_\ell,\varepsilon\rangle,
\quad \Lambda,V\ge 0.
\end{equation}

\paragraph{Primal solution (softmax form).}
If $(\Lambda^\star,V^\star)$ minimizes $F$, then an optimal predictor is
\begin{equation}
\boxed{
\pi_{\Lambda^\star,V^\star}(\ell\mid x)
=
\sigma_\ell\!\left(
\beta\big(\langle \lambda^\star_{\ell'}-\nu^\star_{\ell'},\,t(x)\rangle - r_{\ell'}(x)\big)_{\ell'\in[[L]]}
\right).
}
\end{equation}

\paragraph{Interpretation.}
When $\Lambda=V=0$, (no fairness constraints), we have
\begin{equation}
\pi_{0,0}(\ell\mid x)\propto \exp\big(-\beta(\eta(x)-\widehat y_\ell)^2\big),
\end{equation}
so the distribution concentrates around the grid point closest to $\eta(x)$; the dual terms then
reweight bins to enforce DP across groups.

\subsection{Regularity + why gradient norm matters (Lemma 3.4 + 3.5)}
Let
\begin{equation}
\sigma^2 := 2\sum_{s\in[K]}\frac{1-p_s}{p_s}.
\end{equation}
Then $\nabla F$ is $(\beta\sigma^2)$-Lipschitz, so larger $\beta$ makes $F$ less smooth.

Since the dual domain is constrained ($\Lambda,V\ge 0$), they use the \textbf{gradient mapping}
(for stepsize $\alpha>0$)
\begin{equation}
G_\alpha(\Lambda,V):=
\frac{(\Lambda,V)-\big((\Lambda,V)-\alpha\nabla F(\Lambda,V)\big)_+}{\alpha},
\end{equation}
where $(\cdot)_+$ is coordinatewise projection onto $\R_+^{2(2L+1)K}$.

\paragraph{Lemma 3.5 (key implication).}
For any $(\Lambda,V)\ge 0$, the \emph{fairness violation} of $\pi_{\Lambda,V}$ is controlled by
$\|G_\alpha(\Lambda,V)\|$, and the \emph{risk suboptimality} is bounded by a term linear in
$\|G_\alpha(\Lambda,V)\|$ plus the entropic bias $\log|\widehat{\mathcal Y}_L|/\beta$.

Moral: they need algorithms that directly control $\E\|G_\alpha\|^2$ with clean statistical rates.

\section{Algorithm: online unlabeled post-processing (Algorithm 1)}
Inputs: $L$, $\beta$, number of stochastic gradient calls $T$, $p$, $B$, and functions $\eta,\tau$.
\begin{enumerate}
  \item Build grid $\widehat{\mathcal Y}_L$.
  \item Define $F(\Lambda,V)$ as above; set smoothness proxy $M=\beta\sigma^2$.
  \item Run a black-box stochastic optimizer $\mathcal{A}$ on $F$ for $T$ stochastic gradients to get $(\widehat\Lambda,\widehat V)$.
  \item Output $\pi_{\widehat\Lambda,\widehat V}$ via the softmax formula.
\end{enumerate}

\paragraph{Stochastic gradient oracle uses only $X$.}
They derive $\nabla F$ explicitly and estimate it by sampling a single $X\sim P_X$ and removing the expectation.
So the post-processing step can be done on a \textbf{stream} of unlabeled $X$ without storing data.

\paragraph{Unknown $\eta$ and $\tau$.}
Train $\widehat\eta$ and $\widehat\tau$ on a labeled dataset, then run the same post-processing
on an independent unlabeled stream; additional error terms appear that depend on the quality of
$\widehat\eta,\widehat\tau$.

\section{Main guarantees (Theorem 5.1 + plugin extension)}
They choose parameters as functions of the number of unlabeled samples/iterations $T$
(e.g.\ $L\asymp \sqrt{T}$ and a matching $\beta$ schedule) and prove that:
\begin{itemize}
  \item fairness violation decreases at a rate essentially $1/\sqrt{T}$ (up to logs and constants),
  \item excess risk (vs.\ a natural deterministic DP benchmark based on Kolmogorov--Smirnov-type constraints)
        also decreases like $1/\sqrt{T}$,
  \item dependence on $K$ appears through $\sigma$ and on constraint tightness through $\|(\Lambda^\star,V^\star)\|$.
\end{itemize}
They also provide post-processing-style bounds for the plug-in version with $(\widehat\eta,\widehat\tau)$,
separating estimation error from optimization/post-processing error.

\section{what is actually run}
Datasets:
\begin{itemize}
  \item Law School: predict GPA (scaled to $[0,1]$),
  \item Communities \& Crime: predict violent crime rate (scaled to $[0,1]$),
\end{itemize}
with ethnicity as sensitive attribute (binary: white vs non-white).

Pipeline:
split into train / unlabeled / test (0.4 / 0.4 / 0.2);
train $\widehat\eta$ via LinearRegression and $\widehat\tau$ via LogisticRegression;
run Algorithm 1 using the unlabeled split; evaluate risk and unfairness vs iterations.

They compare against Agarwal--Dudik--Wu (2019)in two training regimes
(with and without access to unlabeled data) $\implies$ competitive or better tradeoffs.

\section{Formules}
\begin{itemize}
  \item Grid indices: $[[L]]=\{-L,\ldots,L\}$, grid values $\widehat y_\ell=\ell B/L$.
  \item Slack DP constraints: $U_s(\pi,\ell)\le \varepsilon_s$.
  \item Dual objective: $F(\Lambda,V)=\E[\mathrm{LSE}_\beta(\cdots)]+\sum_\ell \langle\lambda_\ell+\nu_\ell,\varepsilon\rangle$.
  \item Output predictor: $\pi_{\Lambda,V}(\ell\mid x)=\mathrm{softmax}_\ell\big(\beta(\langle\lambda_{\ell'}-\nu_{\ell'},t(x)\rangle-r_{\ell'}(x))\big)$.
  \item Gradient mapping: $G_\alpha(\Lambda,V)=\frac{(\Lambda,V)-((\Lambda,V)-\alpha\nabla F)_+}{\alpha}$.
\end{itemize}

\newpage 
\section{Comparison: OT characterization vs online convex post-processing}

\subsection*{One-line takeaway}
\begin{itemize}
  \item \textbf{Divol--Gaucher:} population-level characterization of the optimal DP regressor via an OT barycenter
        on a lifted space $(\eta,\Delta)$ (mainly two groups).
  \item \textbf{Taturyan--Chzhen--Hebiri:} practical post-processing algorithm with finite-sample bounds,
        based on discretization + entropic smoothing + stochastic convex optimization, usable with an unlabeled stream of $X$.
\end{itemize}

\subsection*{Differences}
\begin{center}
\begin{tabular}{@{}p{3.2cm}p{6.2cm}p{6.2cm}@{}}
\toprule
 & \textbf{OT} & \textbf{Unlabeled post-processing} \\
\midrule
Primary goal
& Characterize $f^\star$ solving \(\min_f \E[(Y-f(X))^2]\) s.t.\ \(f(X)\perp S\) (population optimum).
& IMPLEMENTABLE estimator achieving (approx.) DP with rates, using post-processing. \\
\addlinespace
Groups
& Focuses on \(\,S\in\{1,2\}\) (extension to $K>2$ left for future work).
& Handles \(\,S\in[K]\) and discusses multiple protected attributes. \\
\addlinespace
Form of predictor
& Deterministic regressor \(f^\star:\X\to\R\), described via OT on a lift \((\eta(x),\Delta(x))\).
& Randomized predictor \(\pi(\cdot\mid x)\) supported on a \textbf{grid} \(\widehat{\mathcal Y}_L\) (can output a sample or the mean). \\
\addlinespace
How DP is enforced
& \textbf{Exact} DP constraint at the population level (\(f(X)\perp S\)).
& \textbf{Relaxed / discretized} DP constraints on grid atoms (with slack \(\varepsilon\)); plus entropic regularization. \\
\addlinespace
Key mathematical object
& OT barycenter with a \emph{multi-to-one} cost after lifting; DP governed by ``disagreement'' regions and a signed density gap \(\Delta\).
& Smooth convex dual objective \(F(\Lambda,V)\); optimal \(\pi\) is an explicit \textbf{softmax} (Lemma 3.1). \\
\addlinespace
What information is implicitly needed
& Needs \(\eta(x)\) and \(\Delta(x)\), where \(\Delta\) depends on group-membership probabilities \(\Pp(S=s\mid X=x)\) (implicit use of $S$).
& Needs \(\eta(x)\) and \(\tau(x)=(\Pp(S=s\mid X=x))_{s}\) to form \(t(x)=\mathbf{1}-\tau(x)/p\). \\
\addlinespace
Computation / “online”
& typically offline
& Stochastic convex optimization with an oracle computable from \emph{unlabeled} i.i.d.\ samples of \(X\); suited for streaming. \\
\bottomrule
\end{tabular}
\end{center}

Fainress through $\Delta(x)$ or $\tau(x)$.

\begin{itemize}
  \item \textbf{OT characterization} more interpretable:
        what the optimal $f^\star$ looks like, why order-preservation can fail, and how DP reshapes predictions geometrically.
  \item \textbf{unlabeled post-processing} more applicable \emph{pipeline}:
        a concrete algorithm with controllable approximation  $(L,\beta,\varepsilon)$ and finite-sample/optimization guarantees.
\end{itemize}

$\implies$ interesting to see the statiqsticak guarantees we can have with the OT post-processing + see what happens with $K>2$.


\clearpage
\part{Comparison on real data}
\label{part:comparison}


\marie{J'ai importé leur code pour utiliser exactement la meme chose, j'ai testé plein de truc mais j'ai toujours des matmul warning de l'enfer meme en scalant tout
par contre j'ai bien verifié et toutes les données sont finies donc c'est harmless en terme de perf mais c'est bizarre}

\paragraph{Data splitting and sample sizes.}
After loading $(X,S,y)$, we follow the same preprocessing choices as the original \textsc{FairReg} evaluation code. 
In particular, for \texttt{adult} and \texttt{lawschool} we optionally rescale the target $y$ to $[-1,1]$ (via a min--max transform) and we subsample the dataset to a fixed size of $N=2000$ points. 
For \texttt{communities}, we keep the full dataset (denote by $N$ its size after loading and optional scaling). 
For each run $i$, we then construct three disjoint splits using two successive calls to \texttt{train\_test\_split}, always \textcolor{red}{stratified by the sensitive attribute} $S$ to preserve group proportions. 
First, we split the full dataset into a training set of size $1-\texttt{TEST\_SIZE}$ and a remaining set of size $\texttt{TEST\_SIZE}$. 
Second, we split this remaining set into an unlabeled set and a test set, using
\[
\texttt{train\_size}=\frac{\texttt{TRAIN\_SIZE}}{1-\texttt{TEST\_SIZE}}.
\]
With the default parameters $\texttt{TRAIN\_SIZE}=0.4$ and $\texttt{TEST\_SIZE}=0.2$, this yields the following proportions of the total sample size $N$:
\begin{align*}
N_{\text{train}} &= (1-\texttt{TEST\_SIZE})\,N = 0.8N,\\
N_{\text{unlab}} &= \texttt{TEST\_SIZE}\cdot\frac{\texttt{TRAIN\_SIZE}}{1-\texttt{TEST\_SIZE}}\,N 
= 0.2\cdot\frac{0.4}{0.8}N = 0.1N,\\
N_{\text{test}}  &= \texttt{TEST\_SIZE}\cdot\Bigl(1-\frac{\texttt{TRAIN\_SIZE}}{1-\texttt{TEST\_SIZE}}\Bigr)\,N
= 0.2\cdot\Bigl(1-\frac{0.4}{0.8}\Bigr)N = 0.1N.
\end{align*}
Concretely, for \texttt{adult} and \texttt{lawschool} where $N=2000$, we obtain $1600$ labeled training points, $200$ unlabeled points, and $200$ test points per run. 
The labeled training split is used to fit the baseline regressor $\eta$ (linear regression) and the sensitive-attribute classifier (logistic regression), while the unlabeled split is used to fit \textsc{FairReg} and to estimate the OT-based correction; all methods are finally evaluated on the held-out test split.

\marie{They use only 2000 points  don't know why but i am trying to find out, maybe it's linked to some hyperparameter tuning}

\paragraph{Hyperparameters.}
For \textsc{FairReg}, we use the default hyperparameters from the original code:
\begin{itemize}
  \item the number of points in the grid is $\sqrt{N_{\text{unlab}}}$ (i.e., $14$ for $N_{\text{unlab}}=200$),
  \item the entropic regularization parameter is $\beta = \tfrac{1}{2}\sqrt{N_{\text{unlab}}}\,\log_2 N_{\text{unlab}}$ (i.e., $\approx 54.1$ for $N_{\text{unlab}}=200$), \marie{la temperature a un enorme impact sur les perf du modele du coup j'ai pas laissé auto sinon c'etait trop nul}
  \item the DP slack parameter is $\varepsilon = 0.01$.
\end{itemize}
\section{Adult dataset}

\paragraph{Adult (Census Income).}
The \texttt{Adult} dataset contains individual-level records from the U.S.\ census, described by demographic and employment-related attributes such as age, education, marital status, occupation, hours worked per week, and capital gain/loss. 
The original prediction task is to infer whether an individual earns more than \$50K per year (binary label \texttt{income}); in the pipeline of Fairreg, this target is preprocessed into a numeric response. 
The dataset also includes sensitive attributes such as \texttt{sex} and \texttt{race}. 
\marie{ y'a aussi des fields de variables categoriques qui ont plein de données manquantes sous forme de "?" qu'on drop au preprocessing}\\

\marie{COMMENTAIRE: c'est un peu bizarre ce qu'ils font dans la pipeline parce que pendant le préprocessing au lieu d'utiliser la target de d'habitude ils utilisent l'age c'est pour ça qu ces exemple est pas hyper relevant (of course l'age est pas biaisé en fonction du genre mais au moins ça test sur des vraies valeurs bien confondues les méthodes...)}

\subsection{Analysis of the dataset}

\begin{figure}[h]
  \centering
    \centering
    \includegraphics[width=\linewidth]{figs/paper/compare_ot_vs_fairreg/adult/full/rf/run_10/ot_diagnostics/ot_scatter_XY_colored_by_S.pdf}
    \caption{Scatter plot of $(X,Y)$ colored by $S$ (Adult dataset; run 10).}
\end{figure}


\subsection{OT}


\begin{figure}[h]
  \centering
    \centering
    \includegraphics[width=\linewidth]{figs/paper/compare_ot_vs_fairreg/adult/full/rf/run_10/ot_diagnostics/transport/ot_ot_matrices_full.pdf}
    \caption{OT plan (Adult dataset; run 10).}
\end{figure}
\begin{figure}[h]
  \centering
    \centering
    \includegraphics[width=0.7\linewidth]{figs/paper/compare_ot_vs_fairreg/adult/full/rf/run_10/ot_diagnostics/ot_agg_plan_eta.pdf}
    \caption{Aggregated OT plan (Adult dataset; run 10).}
    \end{figure}

    \begin{figure}[h]
      \centering
      \begin{minipage}[t]{0.48\linewidth}
        \centering
        \includegraphics[width=\linewidth]{figs/paper/compare_ot_vs_fairreg/adult/full/rf/run_10/ot_diagnostics/ot_eta_delta_colored_by_Y.pdf}
        \small Scatter plot of $(X,Y)$ colored by $\Delta$ (Adult dataset; run 10).
      \end{minipage}\hfill
      \begin{minipage}[t]{0.48\linewidth}
        \centering
        \includegraphics[width=\linewidth]{figs/paper/compare_ot_vs_fairreg/adult/full/rf/run_10/ot_diagnostics/ot_base_vs_fair_scatter.pdf}
        \small Base vs fair scatter (Adult dataset; run 10).
      \end{minipage}
      \caption{Side-by-side: colored-by-$\Delta$ vs base-vs-fair scatter (Adult dataset; run 10).}
    \end{figure}

    \begin{figure}[h]
      \centering
      \begin{minipage}[t]{0.48\linewidth}
        \centering
        \includegraphics[width=\linewidth]{figs/paper/compare_ot_vs_fairreg/adult/full/rf/run_10/ot_diagnostics/ot_barycenter_links.pdf}
        \small Barycenter links.
x      \end{minipage}\hfill
      \begin{minipage}[t]{0.48\linewidth}
        \centering
        \includegraphics[width=\linewidth]{figs/paper/compare_ot_vs_fairreg/adult/full/rf/run_10/ot_diagnostics/ot_eta_delta_transport_lines.pdf}
        \small Transport lines in $(\eta,\Delta)$ space.
      \end{minipage}
      \caption{Barycenter links and transport lines (Adult dataset; run 10).}
    \end{figure}

    \subsection{Hyperparameters tuning}

    \begin{figure}[h]
      \centering
        \centering
        \includegraphics[width=\linewidth]{figs/paper/compare_ot_vs_fairreg/adult/full/rf/run_10/fairreg/fairreg_tuning_curve_L_beta.png}
        \caption{OT tuning curve Adult (run 10).}
    \end{figure}
    % \subsubsection{Since the Fairreg seems to be sensible to hyperparams: tuning}

    % \begin{figure}[h]
    %   \centering
    %     \centering
    %     \includegraphics[width=\linewidth]{figs/paper/compare_ot_vs_fairreg/adult/full/rf/run_10/fairreg_tuning_curve_L_beta.png}
    %     \caption{Fairreg tuning curve Adult (run 10).}
    % \end{figure}
\subsection{Unfairness and MSE}
\begin{table}[t]
  \centering
  \begin{tabular}{ll}
    \toprule
    Quantity & Mean $\pm$ Std \\
    \midrule
    \texttt{mse\_base}        & $0.0161 \pm 0.0005$ \\
    \texttt{mse\_fairreg}     & $0.0162 \pm 0.0005$ \\
    \texttt{mse\_ot}          & $0.0170 \pm 0.0004$ \\
    \texttt{dp\_base\_max}    & $0.1588 \pm 0.0101$ \\
    \texttt{dp\_fairreg\_max} & $0.1498 \pm 0.0100$ \\
    \texttt{dp\_ot\_max}      & $0.0292 \pm 0.0096$ \\
    \bottomrule
  \end{tabular}
  \caption{Summary (mean $\pm$ std over runs) Adult.}
  \label{tab:summary}
\end{table}

\begin{figure} \centering
    \centering
    \includegraphics[width=\linewidth]{figs/paper/compare_ot_vs_fairreg/adult/full/rf/run_10/plots/pred_kde_by_group.png}
    \caption{Prediction distributions by group (Adult dataset; run 10).}
\end{figure}


\subsection{Ordering stability and leakage AUC}


\begin{table}[t]
  \centering
  \begin{tabular}{lrrrr}
  \toprule
  Subset & Spearman $\rho$ & Kendall $\tau_b$ & Mean $|\Delta\mathrm{rank}|$ & Top-5\% Jaccard \\
  \midrule
  overall & 0.447 & 0.347 & 0.220 & 0.408 \\
  S=0 & 0.207 & 0.163 & 0.284 & 0.280 \\
  S=1 & 0.587 & 0.459 & 0.181 & 0.474 \\
  \bottomrule
  \end{tabular}
  \caption[Ordering stability between eta-hat and corrected scores.]{Ordering stability between $\hat\eta$ and corrected scores Adult.}
  \label{tab:ordering}
  \end{table}
  
  \begin{table}[t]
  \centering
  \begin{tabular}{lr}
  \toprule
  Quantity & AUC \\
  \midrule
  $\mathrm{AUC}(S\mid Y)$ & 0.567 \\
  $\mathrm{AUC}(S\mid \hat{\eta})$ & 0.653 \\
  $\mathrm{AUC}(S\mid \tilde{y})$ & 0.480 \\
  \bottomrule
  \end{tabular}
  \caption[Leakage AUC summary.]{Leakage AUC Adult: ability to infer $S$ from a scalar score (values near 0.5 indicate low leakage).}
  \label{tab:auc}
  \end{table}
  

  Conclusion: OT better at reducing leakage and preserving order within groups. But the datase is already almost free from leakage so not really relevant. 

  \section{Communities and Crime dataset}

  \paragraph{Communities and Crime.}
The \texttt{Communities and Crime} dataset is community-level: each example corresponds to a U.S.\ community characterized by a high-dimensional set of socio-economic, demographic, housing, and policing variables (e.g., racial composition percentages, income and poverty indicators, employment structure, household composition, immigration, housing characteristics, and law-enforcement resources). 
The canonical target variable is the violent crime rate per population, \texttt{ViolentCrimesPerPop}. 
In many standard versions of the dataset, features are already normalized to the unit interval, and missing entries can occur (depending on the preprocessing). 
The sensitive attribute is obtained by taking the highest percentage of a racial group in the community (among \texttt{White}, \texttt{Black}, \texttt{Asian}, \texttt{Hispanic}, \texttt{NativeAmerican}, \texttt{Other}) and binarizing it into majority-\texttt{White} vs.\ non-majority-\texttt{White} communities.

  \subsection{Analysis of the dataset}

  Too many columns in X for his dataset so no scatter plot of (X,Y) colored by S. (uncomment to see attempt)
  % too many columns in X 
  % \begin{figure}
  %   \centering
  %     \centering
  %     \includegraphics[width=\linewidth]{figs/paper/compare_ot_vs_fairreg/communities/run_10/ot_diagnostics/ot_scatter_XY_colored_by_S.pdf}
  %     \caption{Scatter plot of $(X,Y)$ colored by $S$ (Crime dataset; run 10).}
  % \end{figure}

  \subsection{OT}
  \begin{figure}[h]
    \centering
      \centering
      \includegraphics[width=\linewidth]{figs/paper/compare_ot_vs_fairreg/communities/full/rf/run_10/ot_diagnostics/transport/ot_ot_matrices_full.pdf}
      \caption{OT plan (Crime dataset; run 10).}
  \end{figure}
  \begin{figure}[h]
    \centering
      \centering
      \includegraphics[width=0.7\linewidth]{figs/paper/compare_ot_vs_fairreg/communities/full/rf/run_10/ot_diagnostics/ot_agg_plan_eta.pdf}
      \caption{Aggregated OT plan (Crime dataset; run 10).}
    \end{figure}
    \begin{figure}[h]
      \centering
      \begin{minipage}[t]{0.48\linewidth}
        \centering
        \includegraphics[width=\linewidth]{figs/paper/compare_ot_vs_fairreg/communities/full/rf/run_10/ot_diagnostics/ot_eta_delta_colored_by_Y.pdf}
        \small Scatter plot of $(X,Y)$ colored by $\Delta$ (Crime dataset; run 10).
      \end{minipage}\hfill
      \begin{minipage}[t]{0.48\linewidth}
        \centering
        \includegraphics[width=\linewidth]{figs/paper/compare_ot_vs_fairreg/communities/full/rf/run_10/ot_diagnostics/ot_base_vs_fair_scatter.pdf}
        \small Base vs fair scatter (Crime dataset; run 10).
      \end{minipage}
      \caption{Side-by-side: colored-by-$\Delta$ vs base-vs-fair scatter (Crime dataset; run 10).}
    \end{figure}
    \begin{figure}[h]
      \centering
      \begin{minipage}[t]{0.48\linewidth}
        \centering
        \includegraphics[width=\linewidth]{figs/paper/compare_ot_vs_fairreg/communities/full/rf/run_10/ot_diagnostics/ot_barycenter_links.pdf}
        \small Barycenter links.
      \end{minipage}\hfill
      \begin{minipage}[t]{0.48\linewidth}
        \centering
        \includegraphics[width=\linewidth]{figs/paper/compare_ot_vs_fairreg/communities/full/rf/run_10/ot_diagnostics/ot_eta_delta_transport_lines.pdf}
        \small Transport lines in $(\eta,\Delta)$ space.
      \end{minipage}
      \caption{Barycenter links and transport lines (Crime dataset; run 10).}
    \end{figure}  
  \subsection{Unfairness and MSE}

  \begin{table}[t]
    \centering
    \begin{tabular}{ll}
      \toprule
      Quantity & Mean $\pm$ Std \\
      \midrule
      \texttt{mse\_base}        & $0.0196 \pm 0.0032$ \\
      \texttt{mse\_fairreg}     & $0.0259 \pm 0.0037$ \\
      \texttt{mse\_ot}          & $0.0408 \pm 0.0059$ \\
      \texttt{dp\_base\_max}    & $0.4326 \pm 0.0395$ \\
      \texttt{dp\_fairreg\_max} & $0.2616 \pm 0.0370$ \\
      \texttt{dp\_ot\_max}      & $0.0981 \pm 0.0369$ \\
      \bottomrule
    \end{tabular}
    \caption{Summary (mean $\pm$ std over runs) Communities and crimes.}
    \label{tab:crime_summary}
  \end{table}
  \begin{figure}
    \centering
      \centering
      \includegraphics[width=\linewidth]{figs/paper/compare_ot_vs_fairreg/communities/full/rf/run_10/plots/pred_kde_by_group.png}
      \caption{Prediction distributions by group (Crime dataset; run 10).}
  \end{figure}  
  \subsection{Ordering stability and leakage AUC}

  \begin{table}[t]
    \centering
    \begin{tabular}{lrrrr}
    \toprule
    Subset & Spearman $\rho$ & Kendall $\tau_b$ & Mean $|\Delta\mathrm{rank}|$ & Top-5\% Jaccard \\
    \midrule
    overall & 0.689 & 0.590 & 0.167 & 0.125 \\
    S=0 & 0.526 & 0.436 & 0.235 & 0.333 \\
    S=1 & 0.908 & 0.813 & 0.076 & 0.273 \\
    \bottomrule
    \end{tabular}
    \caption[Ordering stability between eta-hat and corrected scores Communities and crime.]{Ordering stability between $\hat\eta$ and corrected scores.}
    \label{tab:ordering}
    \end{table}
    
    \begin{table}[t]
    \centering
    \begin{tabular}{lr}
    \toprule
    Quantity & AUC \\
    \midrule
    $\mathrm{AUC}(S\mid Y)$ & 0.168 \\
    $\mathrm{AUC}(S\mid \hat{\eta})$ & 0.096 \\
    $\mathrm{AUC}(S\mid \tilde{y})$ & 0.477 \\
    \bottomrule
    \end{tabular}
    \caption[Leakage AUC summary.]{Leakage AUC Comunities and Crime: ability to infer $S$ from a scalar score (values near 0.5 indicate low leakage).}
    \label{tab:auc}
    \end{table}

    Conclusion: OT better at reducing leakage and preserving order within groups. Here the dataset is very biased so relevant.
    THe reordering is quite strong on S=0 group (much more than on S=1). 
    \marie{Maybe it's because it is actually a merge of multiple communities with different characteristics? }




    \subsection{Hyperparamters tuning}
    \begin{figure}[h]
      \centering
        \centering
        \includegraphics[width=\linewidth]{figs/paper/compare_ot_vs_fairreg/communities/full/rf/run_10/fairreg/fairreg_tuning_curve_L_beta.png}
        \caption{Fairreg tuning curve Communities and Crime (run 10).}
    \end{figure}

    \section{Lawschool dataset}

\paragraph{Law School.}
The \texttt{Law School} dataset contains student-level records related to law school admissions and academic outcomes. 
Covariates include academic indicators such as \texttt{lsat} and undergraduate GPA (\texttt{ugpa}), along with demographic and background variables such as \texttt{race}, \texttt{gender}, age, and family income. Outcome variables include standardized grade-point averages (e.g., first-year GPA \texttt{zfygpa} and overall GPA \texttt{zgpa}) and, in some versions, bar exam passage information (\texttt{bar1}).
The sensitive attribute is the race "1" (white) vs. "0" (non-white).

    \subsection{Analysis of the dataset}
    \begin{figure}
      \centering
        \centering
        \includegraphics[width=\linewidth]{figs/paper/compare_ot_vs_fairreg/lawschool/full/rf/run_10/ot_diagnostics/ot_scatter_XY_colored_by_S.pdf}
        \caption{Scatter plot of $(X,Y)$ colored by $S$ (Lawschool dataset; run 10).}
    \end{figure}
    \subsection{OT}
    \begin{figure}[h]
      \centering
        \centering
        \includegraphics[width=\linewidth]{figs/paper/compare_ot_vs_fairreg/lawschool/full/rf/run_10/ot_diagnostics/transport/ot_ot_matrices_full.pdf}
        \caption{OT plan (Lawschool dataset; run 10).}
    \end{figure}
    \begin{figure}[h]
      \centering
        \centering
        \includegraphics[width=0.7\linewidth]{figs/paper/compare_ot_vs_fairreg/lawschool/full/rf/run_10/ot_diagnostics/ot_agg_plan_eta.pdf}
        \caption{Aggregated OT plan (Lawschool dataset; run 10).}
      \end{figure}
      \begin{figure}[h]
        \centering
        \begin{minipage}[t]{0.48\linewidth}
          \centering
          \includegraphics[width=\linewidth]{figs/paper/compare_ot_vs_fairreg/lawschool/full/rf/run_10/ot_diagnostics/ot_eta_delta_colored_by_Y.pdf}
          \small Scatter plot of $(X,Y)$ colored by $\Delta$ (Lawschool dataset; run 10).
        \end{minipage}\hfill
        \begin{minipage}[t]{0.48\linewidth}
          \centering
          \includegraphics[width=\linewidth]{figs/paper/compare_ot_vs_fairreg/lawschool/full/rf/run_10/ot_diagnostics/ot_base_vs_fair_scatter.pdf}
          \small Base vs fair scatter (Lawschool dataset; run 10).
        \end{minipage}
        \caption{Side-by-side: colored-by-$\Delta$ vs base-vs-f air scatter (Lawschool dataset; run 10).}
      \end{figure}
      \begin{figure}[h]
        \centering
        \begin{minipage}[t]{0.48\linewidth}     
          \centering
          \includegraphics[width=\linewidth]{figs/paper/compare_ot_vs_fairreg/lawschool/full/rf/run_10/ot_diagnostics/ot_barycenter_links.pdf}
          \small Barycenter links.
        \end{minipage}\hfill
        \begin{minipage}[t]{0.48\linewidth}
          \centering
          \includegraphics[width=\linewidth]{figs/paper/compare_ot_vs_fairreg/lawschool/full/rf/run_10/ot_diagnostics/ot_eta_delta_transport_lines.pdf}
          \small Transport lines in $(\eta,\Delta)$ space.
        \end{minipage}
        \caption{Barycenter links and transport lines (Lawschool dataset; run 10).}
      \end{figure}  

      \subsection{Hyperparamters tuning}
      \begin{figure}[h]
        \centering
          \centering
          \includegraphics[width=\linewidth]{figs/paper/compare_ot_vs_fairreg/lawschool/full/rf/run_10/fairreg/fairreg_tuning_curve_L_beta.png}
          \caption{Fairreg tuning curve Lawschool (run 10).}
      \end{figure}

      \subsection{Unfairness and MSE}  

    \begin{table}[t]
      \centering
      \begin{tabular}{ll}
      \toprule
      Quantity & Mean $\pm$ Std \\
      \midrule
      \texttt{mse\_base}        & $0.0093 \pm 0.0004$ \\
      \texttt{mse\_fairreg}     & $0.0095 \pm 0.0004$ \\
      \texttt{mse\_ot}          & $0.0097 \pm 0.0004$ \\
      \texttt{dp\_base\_max}    & $0.2274 \pm 0.0244$ \\
      \texttt{dp\_fairreg\_max} & $0.0969 \pm 0.0202$ \\
      \texttt{dp\_ot\_max}      & $0.0465 \pm 0.0209$ \\
      \bottomrule
      \end{tabular}
      \caption{Summary (mean $\pm$ std over runs) Lawschool.}
      \label{tab:lawschool_summary}
    \end{table}

    \begin{figure}
      \centering
        \centering
        \includegraphics[width=\linewidth]{figs/paper/compare_ot_vs_fairreg/lawschool/full/rf/run_10/plots/pred_kde_by_group.png}
        \caption{Prediction distributions by group (Lawschool dataset; run 10).}
    \end{figure}  
    \subsection{Ordering stability and leakage AUC}

    \begin{table}[t]
      \centering
      \begin{tabular}{lrrrr}
      \toprule
      Subset & Spearman $\rho$ & Kendall $\tau_b$ & Mean $|\Delta\mathrm{rank}|$ & Top-5\% Jaccard \\
      \midrule
      overall & 0.900 & 0.763 & 0.091 & 0.667 \\
      S=0 & 0.865 & 0.677 & 0.121 & 1.000 \\
      S=1 & 0.928 & 0.805 & 0.076 & 0.778 \\
      \bottomrule
      \end{tabular}
      \caption[Ordering stability between eta-hat and corrected scores.]{Ordering stability between $\hat\eta$ and corrected scores Lawchool.}
      \label{tab:ordering}
      \end{table}
      
      \begin{table}[t]
      \centering
      \begin{tabular}{lr}
      \toprule
      Quantity & AUC \\
      \midrule
      $\mathrm{AUC}(S\mid Y)$ & 0.612 \\
      $\mathrm{AUC}(S\mid \hat{\eta})$ & 0.637 \\
      $\mathrm{AUC}(S\mid \tilde{y})$ & 0.513 \\
      \bottomrule
      \end{tabular}
      \caption[Leakage AUC summary.]{Leakage AUC Lawschool: ability to infer $S$ from a scalar score (values near 0.5 indicate low leakage).}
      \label{tab:auc}
      \end{table} 

\clearpage
\part{Relaxed unawareness}
\label{part:relaxed}

\section{Setting}
\begin{itemize}
  \item Let $\mathcal X := \mathbb R \times (0,\infty)$ and write $x=(h,a)\in\mathcal X$
  (with $a$ standing for $|d|$ in the paper).
  \item Let $\mu^+,\mu^-\in\mathcal P(\mathcal X)$ be given.
  \item quadratic cost (same as in the paper)
  \begin{equation*}
  c(x,y) \;:=\; \frac{(h-y)^2}{a},\qquad x=(h,a)\in\mathcal X,\; y\in\mathbb R.
\end{equation*}
  \item For $\alpha\in\mathcal P(\mathcal X)$ and $\nu\in\mathcal P(\mathbb R)$ set
  \begin{equation*}
  OT_c(\alpha,\nu) \;:=\; \inf_{\pi\in\Pi(\alpha,\nu)} \int_{\mathcal X\times\mathbb R} c(x,y)\,d\pi(x,y),
\end{equation*}
  where $\Pi(\alpha,\nu)$ denotes the set of couplings between $\alpha$ and $\nu$.
  \item $W_2$ is the usual 2-Wasserstein distance on $\mathcal P(\mathbb R)$.
\end{itemize}

\medskip

\section{Relaxed unawareness OT problem}

Fix $\lambda>0$. We consider the relaxed problem
\begin{equation}\label{eq:P-lambda}
\inf_{\nu_1,\nu_2\in\mathcal P(\mathbb R)}
\left\{
OT_c(\mu^+,\nu_1)+OT_c(\mu^-,\nu_2)
+\lambda\,W_2^2(\nu_1,\nu_2)
\right\}.
\tag{$P_\lambda$}
\end{equation}

\begin{proposition}[Closed-form relax porblem]\label{prop:relaxed_reduction}
Let $x_1=(h_1,a_1),x_2=(h_2,a_2)\in\mathcal X$. Define the pointwise relaxed cost
\begin{equation}\label{eq:Clambda-def}
C_\lambda(x_1,x_2)
\;:=\;
\inf_{y_1,y_2\in\mathbb R}
\left\{
\frac{(h_1-y_1)^2}{a_1}+\frac{(h_2-y_2)^2}{a_2}+\lambda (y_1-y_2)^2
\right\}.
\end{equation}
Then
\begin{equation}\label{eq:OT-reduced}
\inf (P_\lambda)
\;=\;
\inf_{\pi\in\Pi(\mu^+,\mu^-)} \int_{\mathcal X\times\mathcal X} C_\lambda(x_1,x_2)\,d\pi(x_1,x_2).
\end{equation}
Moreover, the minimizer in \eqref{eq:Clambda-def} is unique and given by
\begin{equation}\label{eq:y-star}
y_1^*(x_1,x_2)=\frac{(1+\lambda a_2)\,h_1+\lambda a_1\,h_2}{1+\lambda(a_1+a_2)},
\qquad
y_2^*(x_1,x_2)=\frac{\lambda a_2\,h_1+(1+\lambda a_1)\,h_2}{1+\lambda(a_1+a_2)},
\end{equation}
and the optimal value has the closed form
\begin{equation}\label{eq:Clambda-closed}
C_\lambda(x_1,x_2)
\;=\;
\frac{\lambda}{1+\lambda(a_1+a_2)}\,(h_1-h_2)^2.
\end{equation}
\end{proposition}

\begin{proof}


\textbf{Step 1: comupling and gluing lemma.}
We expand each term in \eqref{eq:P-lambda} with couplings:
\begin{equation*}
  OT_c(\mu^+,\nu_1)=\inf_{\pi_1\in\Pi(\mu^+,\nu_1)}\int c(x_1,y_1)\,d\pi_1,
\quad
OT_c(\mu^-,\nu_2)=\inf_{\pi_2\in\Pi(\mu^-,\nu_2)}\int c(x_2,y_2)\,d\pi_2,
\end{equation*}
\begin{equation*}
  W_2^2(\nu_1,\nu_2)=\inf_{\gamma\in\Pi(\nu_1,\nu_2)}\int (y_1-y_2)^2\,d\gamma.
\end{equation*}
Hence
\begin{equation*}
  \inf (P_\lambda) = 
  \inf_{\nu_1,\nu_2}\;
  \inf_{\substack{
    \pi_1\in\Pi(\mu^+,\nu_1)\\
    \pi_2\in\Pi(\mu^-,\nu_2)\\
    \gamma\in\Pi(\nu_1,\nu_2)
    }}\left\{
\int c(x_1,y_1)\,d\pi_1 + \int c(x_2,y_2)\,d\pi_2 + \lambda \int (y_1-y_2)^2\,d\gamma
\right\}.
\end{equation*}
% If $(pr_{y_1})_\#\pi_1=\nu_1=(pr_{y_1})_\#\gamma$ and $(pr_{y_2})_\#\gamma=\nu_2=(pr_{y_2})_\#\pi_2$, the triple $(\pi_1,\gamma,\pi_2)$ hascompatible marginals and we can apply the gluing lemma from e.g.,. lemma 5.5 in \textcite{Santambrogio2015}
% By the gluing lemma, for any triple $(\pi_1,\gamma,\pi_2)$ with compatible marginals
% there exists a probability measure $\Pi$ on $\mathcal X\times\mathcal X\times\mathbb R\times\mathbb R$
% whose $(x_1,y_1)$-marginal is $\pi_1$, $(y_1,y_2)$-marginal is $\gamma$, and $(x_2,y_2)$-marginal is $\pi_2$.
% Conversely, any such $\Pi$ defines $\pi_1,\pi_2,\gamma$ by marginalization.
% Therefore
% \begin{equation}\label{eq:Pi-form}
% \inf(P_\lambda)
% =
% \inf_{\Pi\in\mathcal A}
% \int_{\mathcal X^2\times\mathbb R^2}
% \left(
% c(x_1,y_1)+c(x_2,y_2)+\lambda (y_1-y_2)^2
% \right)\,d\Pi,
% \end{equation}
% where $\mathcal A$ is the set of probability measures $\Pi$ on $\mathcal X^2\times\mathbb R^2$
% such that $(\mathrm{pr}_{x_1})_\#\Pi=\mu^+$ and $(\mathrm{pr}_{x_2})_\#\Pi=\mu^-$.



% \paragraph{Gluing.}
% Let $\pi_1\in\Pi(\mu^+,\nu_1)$, $\gamma\in\Pi(\nu_1,\nu_2)$, and $\pi_2\in\Pi(\mu^-,\nu_2)$.
% Equivalently, assume there exist random variables
% \[
% (X_1,Y_1)\sim \pi_1,\qquad (Y_1',Y_2)\sim \gamma,\qquad (X_2,Y_2')\sim \pi_2,
% \]
% such that $Y_1\sim \nu_1$ and $Y_2\sim \nu_2$, and moreover the shared marginals match:
% \[
% \Law(Y_1)=\Law(Y_1')=\nu_1,\qquad \Law(Y_2)=\Law(Y_2')=\nu_2.
% \]
% Then, by the gluing lemma, there exists a quadruple of random variables
% \[
% (X_1,X_2,Y_1,Y_2)\in \mathcal X\times\mathcal X\times\R\times\R
% \]
% with joint law $\Pi:=\Law(X_1,X_2,Y_1,Y_2)$ such that the pairwise marginals coincide with the
% original couplings:
% \[
% \Law(X_1,Y_1)=\pi_1,\qquad \Law(Y_1,Y_2)=\gamma,\qquad \Law(X_2,Y_2)=\pi_2.
% \]
% Conversely, given any joint law $\Pi$ on $(X_1,X_2,Y_1,Y_2)$, one recovers feasible couplings by
% marginalization
% \[
% \pi_1:=\Law(X_1,Y_1),\quad \gamma:=\Law(Y_1,Y_2),\quad \pi_2:=\Law(X_2,Y_2),
% \quad\text{and}\quad
% \nu_1:=\Law(Y_1),\ \nu_2:=\Law(Y_2).
% \]
% Therefore the relaxed objective can be written as a single expectation under $\Pi$:
% \[
% \inf(P_\lambda)
% =
% \inf_{\Pi\in\mathcal A}
% \E_{\Pi}\!\left[c(X_1,Y_1)+c(X_2,Y_2)+\lambda (Y_1-Y_2)^2\right],
% \]
% where $\mathcal A$ denotes the set of probability measures $\Pi$ on
% $\mathcal X\times\mathcal X\times\R\times\R$ such that $\Law(X_1)=\mu^+$ and $\Law(X_2)=\mu^-$.

\paragraph{Gluing lemma.}
Let $\pi_1\in\Pi(\mu^+,\nu_1)$, $\gamma\in\Pi(\nu_1,\nu_2)$, and $\pi_2\in\Pi(\mu^-,\nu_2)$, and assume the
\emph{shared marginals are consistent}, i.e.
\begin{equation*}
  (\mathrm{pr}_{y_1})_\#\pi_1=(\mathrm{pr}_{y_1})_\#\gamma=\nu_1,
\qquad
(\mathrm{pr}_{y_2})_\#\gamma=(\mathrm{pr}_{y_2})_\#\pi_2=\nu_2.
\end{equation*}

One may view $\pi_1,\gamma,\pi_2$ as the laws of random pairs
$(X_1,Y_1)$, $(Y_1,Y_2)$ and $(X_2,Y_2)$ whose overlaps $Y_1$ and $Y_2$ have the same distributions.
By the gluing lemma (see, e.g., Lemma~5.5 in \textcite{Santambrogio2015}), there exists a probability measure
$\Pi$ on $\mathcal X\times\mathcal X\times\R\times\R$ such that
\begin{equation*}
  (\mathrm{pr}_{x_1,y_1})_\#\Pi=\pi_1,\qquad
(\mathrm{pr}_{y_1,y_2})_\#\Pi=\gamma,\qquad
(\mathrm{pr}_{x_2,y_2})_\#\Pi=\pi_2.
\end{equation*}

Conversely, any such $\Pi$ determines feasible couplings by marginalization. Hence the relaxed objective can be
expressed using a single joint law:
\begin{equation}\label{eq:Pi-form}
\inf(P_\lambda)
=
\inf_{\Pi\in\mathcal A}
\int_{\mathcal X^2\times \R^2}
\Big(
c(x_1,y_1)+c(x_2,y_2)+\lambda (y_1-y_2)^2
\Big)\,d\Pi,
\end{equation}
where $\mathcal A$ is the set of probability measures $\Pi$ on $\mathcal X^2\times\R^2$ such that
$(\mathrm{pr}_{x_1})_\#\Pi=\mu^+$ and $(\mathrm{pr}_{x_2})_\#\Pi=\mu^-$.


\paragraph{Step 2: reduce to pointwize optimization.}
Let $\Pi\in\mathcal A$ and let $\pi:=(\mathrm{pr}_{x_1,x_2})_\#\Pi\in\Pi(\mu^+,\mu^-)$ (the $(X_1, X_2)$ marginal).
We can write 
\begin{align*}
  \pi_1(dx_1,dy_1) &= \nu_1(dy_1) = \nu_1(dy_1)K_1(dx_1|y_1),\\
  \pi_2(dx_2,dy_2) &= \nu_2(dy_2) = \nu_2(dy_2)K_2(dx_2|y_2),
\end{align*}
where $K_i(\cdot|y_i)$ are the corresponding conditional distributions of $X_i$ knowing $Y_i= y_i$.
Then 
\begin{equation*}
  \Pi(dx_1,dx_2,dy_1,dy_2)
  = \gamma(dy_1, dy_2)\,K_1(dx_1|y_1)\,K_2(dx_2|y_2).
\end{equation*}
\marie{on tire $(Y_1,Y_2) \sim \gamma$, puis $X_1|Y_1 \sim K_1(\cdot|Y_1)$ et $X_2|Y_2 \sim K_2(\cdot|Y_2)$ à partir de $\pi_i$ indépendamment suivant $y_1, Y_2$.}\\
Thereforen, there exists a Markov kernel (i.e., measurable map) $K_{x_1,x_2}$ on $\mathbb R^2$ such that
\begin{equation*}
  \Pi(dx_1,dx_2,dy_1,dy_2)=\pi(dx_1,dx_2)\,K_{x_1,x_2}(dy_1,dy_2).
\end{equation*}

Write the integrand as
\begin{equation*}
  L(x_1,x_2;y_1,y_2):=
c(x_1,y_1)+c(x_2,y_2)+\lambda (y_1-y_2)^2.
\end{equation*}

Then \eqref{eq:Pi-form} becomes
\[
\int_{\mathcal X^2}\left(\int_{\mathbb R^2} L(x_1,x_2;y_1,y_2)\,dK_{x_1,x_2}(y_1,y_2)\right)\,d\pi(x_1,x_2).
\]
For each fixed $(x_1,x_2)$, the map $(y_1,y_2)\mapsto L(x_1,x_2;y_1,y_2)$ is a strictly convex quadratic.
Hence, for any probability measure $K$ on $\mathbb R^2$,
\[
\int L\,dK \;\ge\; \inf_{(y_1,y_2)\in\mathbb R^2} L(x_1,x_2;y_1,y_2),
\]
with equality achieved by choosing $K$ as a Dirac mass at the unique minimizer $\delta_{(y_1^*, y_2*)}$.
Then for a fixed couple $(x_1, x_2)$ the conditional distribution of $(Y_1,Y_2)$ given $(X_1,X_2)=(x_1,x_2)$ that minimizes the inner integral is the Dirac mass at the unique minimizer of $L(x_1,x_2;\cdot,\cdot)$. 
We have
\begin{equation*}
  C_\lambda(x_1,x_2)
  :=
  \inf_{y_1,y_2\in\mathbb R}
  L(x_1,x_2;y_1,y_2)
\end{equation*}
Therefore, for a fixed $\pi\in\Pi(\mu^+,\mu^-)$, the optimization over $\Pi\in\mathcal A$ with this marginal reduces tp an optimization over the kernels $\{K_{x_1,x_2}\}$, and we have 
% \begin{equation*}
%   \int_{\mathcal X^2}\underbrace{\left(\int_{\mathbb R^2} L(x_1,x_2;\;\cdot \;)\,dK_{x_1,x_2}(y_1,y_2)\right)}\,d\pi(x_1,x_2)
% \end{equation*}

\begin{equation*}
  \int_{\mathcal X^2}
  \left(
  \tikzmarknode[fill=targetGreen!75,
               rounded corners=0.4ex,
               inner xsep=0.7ex,
               inner ysep=0.3ex]{innerK}
  {\displaystyle \int_{\mathbb R^2} L(x_1,x_2;\,\cdot\,)\,dK_{x_1,x_2}(y_1,y_2)}
  \right)
  \,d\pi(x_1,x_2).
  \end{equation*}
  
  \begin{tikzpicture}[overlay,remember picture,
    >=Stealth,
    every node/.style={font=\small}]
    \node[below=1.4em of innerK.south, anchor=south, text=targetGreen!70!black] (labK)
      {depends only on $K_{x_1,x_2}$};
  \end{tikzpicture}

  Thus we can minimize the inner integral pointwise. At a fixed $\pi\in\Pi(\mu^+,\mu^-)$ we have that the best kernel choise is 
  \begin{equation*}
    K_{x_1,x_2} = \delta_{(y_1^*(x_1,x_2) \;, \;y_2^*(x_1,x_2))},
  \end{equation*}
  where $(y_1^*(x_1,x_2),y_2^*(x_1,x_2))$ is the unique minimizer of $L(x_1,x_2;\cdot,\cdot)$.
We went from an optimization over $\Pi\in\mathcal A$ with fixed $(\mathrm{pr}_{x_1,x_2})_\#\Pi=\pi$ to a pointwise minimization for each $\pi$ where we know the optimal kernel. 
Hence
\begin{equation*}
  \inf_{\Pi\in\mathcal A:\;(\mathrm{pr}_{x_1,x_2})_\#\Pi=\pi}\;
  \int L\,d\Pi
  \int_{\mathcal X^2}
  \underbrace{\left(
    % \inf_{y_1,y_2\in\mathbb R}
    L(x_1,x_2;y_1,y_2)
    \right)}_{C_\lambda(x_1,x_2)}
  \,d\pi(x_1,x_2).
\end{equation*}
  
where $C_\lambda$ is defined in \eqref{eq:Clambda-def}.
Taking the infimum over $\pi\in\Pi(\mu^+,\mu^-)$ yields \eqref{eq:OT-reduced}.

\paragraph{Step 3: closed form.}
Fix $x_i=(h_i,a_i)$ and minimize
\begin{equation*}
  \Phi(y_1,y_2)=\frac{(h_1-y_1)^2}{a_1}+\frac{(h_2-y_2)^2}{a_2}+\lambda (y_1-y_2)^2.
\end{equation*}

First-order optimality gives
\[ 
\begin{cases}
  \frac{2(y_1-h_1)}{a_1}+2\lambda(y_1-y_2)&=0,\\
\frac{2(y_2-h_2)}{a_2}-2\lambda(y_1-y_2)&=0.
\end{cases}
\]

Solving the system yields \eqref{eq:y-star}.
Let $u:=y_1^*-y_2^*$. From the optimality conditions we also have
\[
\begin{cases}
  y_1^*-h_1&=-\lambda a_1 u,\\
  y_2^*-h_2&=\lambda a_2 u,
\end{cases}
\]
and subtracting these identities gives
\begin{equation*}
  u=(y_1^*-y_2^*)=(h_1-h_2)-\lambda(a_1+a_2)u,
\end{equation*}

and then
\begin{equation*}
  u=\frac{h_1-h_2}{1+\lambda(a_1+a_2)}.
\end{equation*}


Plugging into $\Phi$ at optimum:
\begin{align*}
  C_\lambda(x_1,x_2)
&=\Phi(y_1^*,y_2^*), \\
&=
\frac{(\lambda a_1 u)^2}{a_1}+\frac{(\lambda a_2 u)^2}{a_2}+\lambda u^2, \\
&=
\lambda u^2\big(1+\lambda(a_1+a_2)\big),
\end{align*}

hence
\begin{align*}
  C_\lambda(x_1,x_2)
&=
\lambda \frac{(h_1-h_2)^2}{(1+\lambda(a_1+a_2))^2}\,(1+\lambda(a_1+a_2)), \\
&=
\frac{\lambda}{1+\lambda(a_1+a_2)}(h_1-h_2)^2,
\end{align*}

which is \eqref{eq:Clambda-closed}.
\end{proof}

% \begin{remark}[Reconstruction of $(\nu_1^*,\nu_2^*)$ from an optimal coupling]
% Let $\pi^*\in\Pi(\mu^+,\mu^-)$ be optimal for \eqref{eq:OT-reduced} and define the maps
% $T_1(x_1,x_2):=y_1^*(x_1,x_2)$ and $T_2(x_1,x_2):=y_2^*(x_1,x_2)$ from \eqref{eq:y-star}.
% Then an optimal solution of \eqref{eq:P-lambda} is obtained by
% \[
% \nu_1^* = (T_1)_\#\pi^*,
% \qquad
% \nu_2^* = (T_2)_\#\pi^*,
% \]
% and the optimal $W_2$-coupling between $\nu_1^*$ and $\nu_2^*$ is simply the pushforward
% of $\pi^*$ by $(T_1,T_2)$.
% \end{remark}

\begin{remark}[Sanity check for compatiblity wih he unrelaxed solution]\label{rem:lambda_to_infty}
  Consider the relaxed objective \eqref{eq:P-lambda} and its reduced OT form \eqref{eq:OT-reduced} with
  \begin{equation*}
      C_\lambda(x_1,x_2)=\inf_{y_1,y_2}\Big\{c(x_1,y_1)+c(x_2,y_2)+\lambda(y_1-y_2)^2\Big\}.
    \end{equation*}

  Let $x_i=(h_i,a_i)$ with $a_i>0$. The inner minimization admits the explicit minimizer
  $(y_1^*(x_1,x_2),y_2^*(x_1,x_2))$ given in \eqref{eq:y-star}, and one checks that
  \begin{equation}\label{eq:collapse_y1y2}
  y_1^*(x_1,x_2)-y_2^*(x_1,x_2)=\frac{h_1-h_2}{1+\lambda(a_1+a_2)}\xrightarrow[\lambda\to\infty]{}0,
  \end{equation}
  and moreover
  \begin{equation}\label{eq:collapse_to_bary}
  y_1^*(x_1,x_2),\ y_2^*(x_1,x_2)
  \xrightarrow[\lambda\to\infty]{}
  \frac{a_2h_1+a_1h_2}{a_1+a_2}
  =:y^\star(x_1,x_2),
  \end{equation}
  where $y^\star(x_1,x_2)$ is exactly the minimizer of the unrelaxed two-to-one collapse
  $\inf_{y}\{c(x_1,y)+c(x_2,y)\}$.
  
  At the level of pairwise costs, the closed form \eqref{eq:Clambda-closed} yields the pointwise convergence
  \begin{equation*}
  C_\lambda(x_1,x_2)=\frac{\lambda}{1+\lambda(a_1+a_2)}(h_1-h_2)^2
  \;\xrightarrow[\lambda\to\infty]{}\;
  \frac{(h_1-h_2)^2}{a_1+a_2}
  =:C_\infty(x_1,x_2),
\end{equation*}

  which is precisely the unrelaxed OT cost. In particular, if $V_\lambda$ denotes the optimal value of
  \eqref{eq:OT-reduced}, then $V_\lambda$ is nondecreasing in $\lambda$ and satisfies
  $V_\lambda\uparrow V_\infty$ as $\lambda\to\infty$, where $V_\infty$ is the optimal value of the
  unrelaxed problem.
  
  Finally, let $\pi_\lambda^*\in\Pi(\mu^+,\mu^-)$ be an optimal coupling for the reduced problem
\eqref{eq:OT-reduced}, and define the associated optimal outputs by
\begin{equation*}
  Y_{1,\lambda}:=T_{1,\lambda}(X_1,X_2),\qquad Y_{2,\lambda}:=T_{2,\lambda}(X_1,X_2),
\qquad (X_1,X_2)\sim \pi_\lambda^*,
\end{equation*}

so that $\nu_{1,\lambda}^*=\Law(Y_{1,\lambda})=(T_{1,\lambda})_\#\pi_\lambda^*$ and
$\nu_{2,\lambda}^*=\Law(Y_{2,\lambda})=(T_{2,\lambda})_\#\pi_\lambda^*$.
Consider moreover the coupling between these two output laws induced by the same construction,
\begin{equation*}
  \gamma_\lambda:=(T_{1,\lambda},T_{2,\lambda})_\#\pi_\lambda^*\in\Pi(\nu_{1,\lambda}^*,\nu_{2,\lambda}^*).
\end{equation*}

By definition of $W_2$ we have
\begin{align*}
  W_2^2(\nu_{1,\lambda}^*,\nu_{2,\lambda}^*)
&\;\le\;
\int_{\R^2} (y_1-y_2)^2\,d\gamma_\lambda(y_1,y_2), \\
&\;=\;
\E_{\pi_\lambda^*}\!\big[(Y_{1,\lambda}-Y_{2,\lambda})^2\big].
\end{align*}

Using the explicit formula \eqref{eq:collapse_y1y2}, we obtain that
$Y_{1,\lambda}-Y_{2,\lambda}\to 0$ pointwise in $(x_1,x_2)$ as $\lambda\to\infty$, and hence (under the
usual integrability ensured here by the quadratic structure) the right-hand side vanishes as $\lambda\to\infty$.
In particular, the two output laws become arbitrarily close in $W_2$:
\begin{equation*}
  W_2(\nu_{1,\lambda}^*,\nu_{2,\lambda}^*)\xrightarrow[\lambda\to\infty]{}0.
\end{equation*}

Therefore the relaxed solutions asymptotically enforce a single common output distribution; if the unrelaxed
barycenter is unique, this common limit must coincide with the unrelaxed barycenter law.

  
  \end{remark}


\newpage
\printbibliography





\end{document}
